% A CPI stack (Lesson 8): the CPI of 1.31 of the five-stage example
% split into its base (steady-state) component of 1.00 and the
% contributions of branch mispredictions (0.06) and cache misses (0.25).
% Original.
% Usage: \resizebox{\linewidth}{!}{% A CPI stack (Lesson 8): the CPI of 1.31 of the five-stage example
% split into its base (steady-state) component of 1.00 and the
% contributions of branch mispredictions (0.06) and cache misses (0.25).
% Original.
% Usage: \resizebox{\linewidth}{!}{% A CPI stack (Lesson 8): the CPI of 1.31 of the five-stage example
% split into its base (steady-state) component of 1.00 and the
% contributions of branch mispredictions (0.06) and cache misses (0.25).
% Original.
% Usage: \resizebox{\linewidth}{!}{% A CPI stack (Lesson 8): the CPI of 1.31 of the five-stage example
% split into its base (steady-state) component of 1.00 and the
% contributions of branch mispredictions (0.06) and cache misses (0.25).
% Original.
% Usage: \resizebox{\linewidth}{!}{\input{figures/l08-cpi-stack}}
% Width: margin (50 mm); the drawing is about 45 mm wide.
\begin{tikzpicture}[x=1mm, y=24mm, font=\scriptsize\sffamily]
  \draw[fill=ln-accent!20] (0,0) rectangle (12,1.00)
    node[midway, align=center] {\iflangja{基本}{base}\\1.00};
  \draw[fill=ln-tip!30] (0,1.00) rectangle (12,1.06);
  \draw[fill=ln-caution!25] (0,1.06) rectangle (12,1.31);
  \draw[ln-muted] (12,1.03) -- (15,1.03)
    node[right, text=black] {\iflangja{分岐予測ミス 0.06}{branch mispred.\ 0.06}};
  \draw[ln-muted] (12,1.185) -- (15,1.185)
    node[right, text=black] {\iflangja{キャッシュミス 0.25}{cache misses 0.25}};
  \draw[->] (-2,0) -- (-2,1.45) node[above, inner sep=1pt] {CPI};
  \foreach \v in {0,0.5,1} {
    \draw (-2,\v) -- (-2.8,\v) node[left, inner sep=1pt] {\v};
  }
  \draw[dashed, ln-muted] (-2,1.31) -- (0,1.31);
  \node[left, inner sep=1pt] at (-2.8,1.31) {1.31};
\end{tikzpicture}
}
% Width: margin (50 mm); the drawing is about 45 mm wide.
\begin{tikzpicture}[x=1mm, y=24mm, font=\scriptsize\sffamily]
  \draw[fill=ln-accent!20] (0,0) rectangle (12,1.00)
    node[midway, align=center] {\iflangja{基本}{base}\\1.00};
  \draw[fill=ln-tip!30] (0,1.00) rectangle (12,1.06);
  \draw[fill=ln-caution!25] (0,1.06) rectangle (12,1.31);
  \draw[ln-muted] (12,1.03) -- (15,1.03)
    node[right, text=black] {\iflangja{分岐予測ミス 0.06}{branch mispred.\ 0.06}};
  \draw[ln-muted] (12,1.185) -- (15,1.185)
    node[right, text=black] {\iflangja{キャッシュミス 0.25}{cache misses 0.25}};
  \draw[->] (-2,0) -- (-2,1.45) node[above, inner sep=1pt] {CPI};
  \foreach \v in {0,0.5,1} {
    \draw (-2,\v) -- (-2.8,\v) node[left, inner sep=1pt] {\v};
  }
  \draw[dashed, ln-muted] (-2,1.31) -- (0,1.31);
  \node[left, inner sep=1pt] at (-2.8,1.31) {1.31};
\end{tikzpicture}
}
% Width: margin (50 mm); the drawing is about 45 mm wide.
\begin{tikzpicture}[x=1mm, y=24mm, font=\scriptsize\sffamily]
  \draw[fill=ln-accent!20] (0,0) rectangle (12,1.00)
    node[midway, align=center] {\iflangja{基本}{base}\\1.00};
  \draw[fill=ln-tip!30] (0,1.00) rectangle (12,1.06);
  \draw[fill=ln-caution!25] (0,1.06) rectangle (12,1.31);
  \draw[ln-muted] (12,1.03) -- (15,1.03)
    node[right, text=black] {\iflangja{分岐予測ミス 0.06}{branch mispred.\ 0.06}};
  \draw[ln-muted] (12,1.185) -- (15,1.185)
    node[right, text=black] {\iflangja{キャッシュミス 0.25}{cache misses 0.25}};
  \draw[->] (-2,0) -- (-2,1.45) node[above, inner sep=1pt] {CPI};
  \foreach \v in {0,0.5,1} {
    \draw (-2,\v) -- (-2.8,\v) node[left, inner sep=1pt] {\v};
  }
  \draw[dashed, ln-muted] (-2,1.31) -- (0,1.31);
  \node[left, inner sep=1pt] at (-2.8,1.31) {1.31};
\end{tikzpicture}
}
% Width: margin (50 mm); the drawing is about 45 mm wide.
\begin{tikzpicture}[x=1mm, y=24mm, font=\scriptsize\sffamily]
  \draw[fill=ln-accent!20] (0,0) rectangle (12,1.00)
    node[midway, align=center] {\iflangja{基本}{base}\\1.00};
  \draw[fill=ln-tip!30] (0,1.00) rectangle (12,1.06);
  \draw[fill=ln-caution!25] (0,1.06) rectangle (12,1.31);
  \draw[ln-muted] (12,1.03) -- (15,1.03)
    node[right, text=black] {\iflangja{分岐予測ミス 0.06}{branch mispred.\ 0.06}};
  \draw[ln-muted] (12,1.185) -- (15,1.185)
    node[right, text=black] {\iflangja{キャッシュミス 0.25}{cache misses 0.25}};
  \draw[->] (-2,0) -- (-2,1.45) node[above, inner sep=1pt] {CPI};
  \foreach \v in {0,0.5,1} {
    \draw (-2,\v) -- (-2.8,\v) node[left, inner sep=1pt] {\v};
  }
  \draw[dashed, ln-muted] (-2,1.31) -- (0,1.31);
  \node[left, inner sep=1pt] at (-2.8,1.31) {1.31};
\end{tikzpicture}
